% !tp-example Basic usage
\documentclass{article}
\begin{document}
\rawhtml|<tp-multi-choice-question answer="1,3,5" random="">
  <dl>
    <dt>Title</dt><dd>Web languages</dd>
    <dt>Prompt</dt><dd>Select the styling languages, including CSS preprocessors.</dd>
    <dt>Form</dt><dd><ul><li>CSS</li><li>HTML</li><li>Sass</li><li>JavaScript</li><li>Less</li></ul></dd>
    <dt>Feedback</dt><dd><ul><li>Correct. CSS describes presentation.</li><li>HTML structures content.</li><li>Correct. Sass is a CSS preprocessor.</li><li>JavaScript is a programming language.</li><li>Correct. Less is a CSS preprocessor.</li></ul></dd>
    <dt>Solution</dt><dd>CSS, Sass and Less are used to define styles.</dd>
  </dl>
</tp-multi-choice-question>|
\end{document}

% !tp-example Javascript API
\documentclass{article}
\begin{document}
\wcScript[type=module]|
import "/components/multi-choice-question/multi-choice-question.js";

const example = document.createDocumentFragment();

const element1 = document.createElement("tp-multi-choice-question");
element1.setAttribute("answer", "1,3,5");
element1.setAttribute("random", "");
element1.append(document.createTextNode("\n  "));

const element2 = document.createElement("dl");
element2.append(document.createTextNode("\n    "));

const element3 = document.createElement("dt");
element3.append(document.createTextNode("Title"));
element2.append(element3);

const element4 = document.createElement("dd");
element4.append(document.createTextNode("Web languages"));
element2.append(element4);
element2.append(document.createTextNode("\n    "));

const element5 = document.createElement("dt");
element5.append(document.createTextNode("Prompt"));
element2.append(element5);

const element6 = document.createElement("dd");
element6.append(document.createTextNode("Select the styling languages, including CSS preprocessors."));
element2.append(element6);
element2.append(document.createTextNode("\n    "));

const element7 = document.createElement("dt");
element7.append(document.createTextNode("Form"));
element2.append(element7);

const element8 = document.createElement("dd");

const element9 = document.createElement("ul");

const element10 = document.createElement("li");
element10.append(document.createTextNode("CSS"));
element9.append(element10);

const element11 = document.createElement("li");
element11.append(document.createTextNode("HTML"));
element9.append(element11);

const element12 = document.createElement("li");
element12.append(document.createTextNode("Sass"));
element9.append(element12);

const element13 = document.createElement("li");
element13.append(document.createTextNode("JavaScript"));
element9.append(element13);

const element14 = document.createElement("li");
element14.append(document.createTextNode("Less"));
element9.append(element14);
element8.append(element9);
element2.append(element8);
element2.append(document.createTextNode("\n    "));

const element15 = document.createElement("dt");
element15.append(document.createTextNode("Feedback"));
element2.append(element15);

const element16 = document.createElement("dd");

const element17 = document.createElement("ul");

const element18 = document.createElement("li");
element18.append(document.createTextNode("Correct. CSS describes presentation."));
element17.append(element18);

const element19 = document.createElement("li");
element19.append(document.createTextNode("HTML structures content."));
element17.append(element19);

const element20 = document.createElement("li");
element20.append(document.createTextNode("Correct. Sass is a CSS preprocessor."));
element17.append(element20);

const element21 = document.createElement("li");
element21.append(document.createTextNode("JavaScript is a programming language."));
element17.append(element21);

const element22 = document.createElement("li");
element22.append(document.createTextNode("Correct. Less is a CSS preprocessor."));
element17.append(element22);
element16.append(element17);
element2.append(element16);
element2.append(document.createTextNode("\n    "));

const element23 = document.createElement("dt");
element23.append(document.createTextNode("Solution"));
element2.append(element23);

const element24 = document.createElement("dd");
element24.append(document.createTextNode("CSS, Sass and Less are used to define styles."));
element2.append(element24);
element2.append(document.createTextNode("\n  "));
element1.append(element2);
element1.append(document.createTextNode("\n"));
example.append(element1);

document.body.append(example);
|
\end{document}

% !tp-example Derivatives of tan(x)
\documentclass{article}
\begin{document}
\rawhtml|<tp-multi-choice-question answer="1,3" lang="en" random="">
  <dl>
    <dt>Title</dt><dd>Equivalent derivatives of the tangent</dd>
    <dt>Prompt</dt><dd>Select all expressions equal to the derivative of \(f(x) = \tan(x)\) wherever \(\cos(x) \ne 0\).</dd>
    <dt>Form</dt><dd><ul><li>\(\displaystyle \frac{1}{\cos^2(x)}\)</li><li>\(1 - \tan^2(x)\)</li><li>\(1 + \tan^2(x)\)</li><li>\(\cos^2(x)\)</li><li>\(\displaystyle -\frac{1}{\sin^2(x)}\)</li></ul></dd>
    <dt>Feedback</dt><dd><ul>
      <li>Correct. Apply the quotient rule to sine divided by cosine.</li>
      <li>The identity uses a plus sign, not a minus sign.</li>
      <li>Correct. This is equivalent to the reciprocal square of the cosine.</li>
      <li>The square of the cosine belongs in the denominator.</li>
      <li>This is the derivative of the cotangent, not the tangent.</li>
    </ul></dd>
    <dt>Solution</dt><dd>\(\displaystyle f'(x) = \frac{1}{\cos^2(x)} = 1 + \tan^2(x)\). The expression \(1 - \tan^2(x)\) is not the derivative.</dd>
  </dl>
</tp-multi-choice-question>|

\wcScript|
  // Typeset after the question has created its checkbox controls.
  window.MathJax = {
    svg: { fontCache: "local" },
    startup: {
      pageReady: async () => {
        await customElements.whenDefined("tp-multi-choice-question");
        await customElements.whenDefined("tp-checkbox-list");
        return MathJax.startup.defaultPageReady();
      }
    }
  };
|

\wcScript[defer,src=https://cdn.jsdelivr.net/npm/mathjax@4/tex-svg.js]||
\end{document}

% !tp-example Markup languages
\documentclass{article}
\begin{document}
\rawhtml|<tp-multi-choice-question answer="1,3,5" lang="en" orientation="horizontal" random="">
  <dl>
    <dt>Title</dt><dd>Recognize markup languages</dd>
    <dt>Prompt</dt><dd>Select the three logos representing markup languages rather than programming languages.</dd>
    <dt>Form</dt><dd><ul><li><tp-icon name="file_type_html" library="languages" size="2em" role="img" aria-label="HTML"></tp-icon></li><li><tp-icon name="file_type_javascript" library="languages" size="2em" role="img" aria-label="JavaScript"></tp-icon></li><li><tp-icon name="file_type_markdown" library="languages" size="2em" role="img" aria-label="Markdown"></tp-icon></li><li><tp-icon name="file_type_python" library="languages" size="2em" role="img" aria-label="Python"></tp-icon></li><li><tp-icon name="file_type_asciidoc" library="languages" size="2em" role="img" aria-label="AsciiDoc"></tp-icon></li><li><tp-icon name="file_type_typescript" library="languages" size="2em" role="img" aria-label="TypeScript"></tp-icon></li></ul></dd>
    <dt>Feedback</dt><dd><ul>
      <li>Correct. HTML marks up the structure of web documents.</li>
      <li>JavaScript is a programming language.</li>
      <li>Correct. Markdown is a lightweight markup language.</li>
      <li>Python is a programming language.</li>
      <li>Correct. AsciiDoc is a markup language for documents.</li>
      <li>TypeScript is a programming language based on JavaScript.</li>
    </ul></dd>
    <dt>Solution</dt><dd><p><tp-icon name="file_type_html" library="languages" size="1em" aria-hidden="true"></tp-icon> HTML, <tp-icon name="file_type_markdown" library="languages" size="1em" aria-hidden="true"></tp-icon> Markdown and <tp-icon name="file_type_asciidoc" library="languages" size="1em" aria-hidden="true"></tp-icon> AsciiDoc are markup languages.</p><p><tp-icon name="file_type_javascript" library="languages" size="1em" aria-hidden="true"></tp-icon> JavaScript, <tp-icon name="file_type_python" library="languages" size="1em" aria-hidden="true"></tp-icon> Python and <tp-icon name="file_type_typescript" library="languages" size="1em" aria-hidden="true"></tp-icon> TypeScript are programming languages.</p></dd>
  </dl>
</tp-multi-choice-question>|
\end{document}

% !tp-example No correct choices
\documentclass{article}
\begin{document}
\wcTpMultiChoiceQuestion[answer]{
  \begin{description}
    \item[{Title}] No correct choices
    \item[{Prompt}] Select the even numbers. If none applies, submit without selecting a box.
    \item[{Form}] \begin{itemize}
      \item 3
      \item 5
      \item 7
    \end{itemize}
    \item[{Solution}] All three numbers are odd. Leaving every box unchecked is correct.
  \end{description}
}
\end{document}
